% GENERATED FILE. Edit canonical lessons, then run npm run generate:books.
% canonical-source-hash: fnv1a64:a202182f7b972038
% canonical-lessons: TE-C146-ceru, TE-C146-vesuko, TE-C146-pilucu, TE-C146-lekkapettu, TE-C146-kadugu

\chapter{To reach, To put on, To call, To count, To wash}
\label{ch:te-to-reach-to-put-on-to-call-to-count-to-wash}
\hlchaptermodality{146}

\begin{chapteropening}
\textbf{By the end of this chapter:} \emph{I can say to reach, to put on, to call, to count and to wash.}

Complete the last lesson of chapter 146: i can say to reach, to put on, to call, to count and to wash.
\end{chapteropening}

\section[\texorpdfstring{\te{చేరు} (cēru)}{చేరు (cēru)}]{\te{చేరు} (cēru) — to reach, to join}

\label{lesson:TE-C146-ceru}



\begin{quote}
Before the new one: say the Telugu for to jump, then the Telugu for to get down.
\end{quote}

\subsection*{You'll want to know: \te{చేరు}}
\textbf{\te{చేరు}} — \emph{cēru} — \textquotedblleft{}to reach, to join\textquotedblright{}.

\begin{grammarlens}[title={What you've built}]
One more word for this chapter.
\end{grammarlens}

\subsection*{Your turn}
\begin{itemize}
\raggedright
  \item \emph{Say it:} \emph{cēru}
  \item \emph{Say it:} \emph{cēru}, once more
  \item \emph{Say it:} say \emph{cēru}
  \item \emph{Recall:} say the Telugu for to jump, then the Telugu for to get down, then say \emph{cēru} again
\end{itemize}

\begin{tcolorbox}[breakable,colback=teal!4,colframe=teal!35!black,title={Before you move on}]
What does \te{చేరు} mean? (\textquotedblleft{}To reach, to join\textquotedblright{}.) Say it once more.
\end{tcolorbox}
\section[\texorpdfstring{\te{వేసుకో} (vēsukō)}{వేసుకో (vēsukō)}]{\te{వేసుకో} (vēsukō) — to put on (clothes)}

\label{lesson:TE-C146-vesuko}



\begin{quote}
Before the new one: say the Telugu for to get down, then the Telugu for to reach.
\end{quote}

\subsection*{You'll want to know: \te{వేసుకో}}
\textbf{\te{వేసుకో}} — \emph{vēsukō} — \textquotedblleft{}to put on (clothes)\textquotedblright{}.

\begin{grammarlens}[title={What you've built}]
One more word for this chapter.
\end{grammarlens}

\subsection*{Your turn}
\begin{itemize}
\raggedright
  \item \emph{Say it:} \emph{vēsukō}
  \item \emph{Say it:} \emph{vēsukō}, once more
  \item \emph{Say it:} say \emph{vēsukō}
  \item \emph{Recall:} say the Telugu for to get down, then the Telugu for to reach, then say \emph{vēsukō} again
\end{itemize}

\begin{tcolorbox}[breakable,colback=teal!4,colframe=teal!35!black,title={Before you move on}]
What does \te{వేసుకో} mean? (\textquotedblleft{}To put on (clothes)\textquotedblright{}.) Say it once more.
\end{tcolorbox}
\section[\texorpdfstring{\te{పిలుచు} (pilucu)}{పిలుచు (pilucu)}]{\te{పిలుచు} (pilucu) — to call}

\label{lesson:TE-C146-pilucu}



\begin{quote}
Before the new one: say the Telugu for to reach, then the Telugu for to put on.
\end{quote}

\subsection*{You'll want to know: \te{పిలుచు}}
\textbf{\te{పిలుచు}} — \emph{pilucu} — \textquotedblleft{}to call\textquotedblright{}.

\begin{grammarlens}[title={What you've built}]
One more word for this chapter.
\end{grammarlens}

\subsection*{Your turn}
\begin{itemize}
\raggedright
  \item \emph{Say it:} \emph{pilucu}
  \item \emph{Say it:} \emph{pilucu}, once more
  \item \emph{Say it:} say \emph{pilucu}
  \item \emph{Recall:} say the Telugu for to reach, then the Telugu for to put on, then say \emph{pilucu} again
\end{itemize}

\begin{tcolorbox}[breakable,colback=teal!4,colframe=teal!35!black,title={Before you move on}]
What does \te{పిలుచు} mean? (\textquotedblleft{}To call\textquotedblright{}.) Say it once more.
\end{tcolorbox}
\section[\texorpdfstring{\te{లెక్కపెట్టు} (lekkapeṭṭu)}{లెక్కపెట్టు (lekkapeṭṭu)}]{\te{లెక్కపెట్టు} (lekkapeṭṭu) — to count}

\label{lesson:TE-C146-lekkapettu}



\begin{quote}
Before the new one: say the Telugu for to put on, then the Telugu for to call.

Count on with \textbf{-\te{వ}}: \textbf{\te{పదవ}}, then \textbf{\te{పదకొండవ}}.
\end{quote}

\subsection*{You'll want to know: \te{లెక్కపెట్టు}}
\textbf{\te{లెక్కపెట్టు}} — \emph{lekkapeṭṭu} — \textquotedblleft{}to count\textquotedblright{}.

\begin{grammarlens}[title={What you've built}]
One more word for this chapter.
\end{grammarlens}

\subsection*{Your turn}
\begin{itemize}
\raggedright
  \item \emph{Say it:} \emph{lekkapeṭṭu}
  \item \emph{Say it:} \emph{lekkapeṭṭu}, once more
  \item \emph{Say it:} say \emph{lekkapeṭṭu}
  \item \emph{Recall:} say the Telugu for to put on, then the Telugu for to call, then say \emph{lekkapeṭṭu} again
\end{itemize}

\begin{tcolorbox}[breakable,colback=teal!4,colframe=teal!35!black,title={Before you move on}]
What does \te{లెక్కపెట్టు} mean? (\textquotedblleft{}To count\textquotedblright{}.) Say it once more.
\end{tcolorbox}
\section[\texorpdfstring{\te{కడుగు} (kaḍugu)}{కడుగు (kaḍugu)}]{\te{కడుగు} (kaḍugu) — to wash}

\label{lesson:TE-C146-kadugu}



\begin{quote}
Before the new one: say the Telugu for to call, then the Telugu for to count.
\end{quote}

\subsection*{You'll want to know: \te{కడుగు}}
\textbf{\te{కడుగు}} — \emph{kaḍugu} — \textquotedblleft{}to wash\textquotedblright{}.

\begin{grammarlens}[title={What you've built}]
One more word for this chapter.
\end{grammarlens}

\subsection*{Your turn}
\begin{itemize}
\raggedright
  \item \emph{Say it:} \emph{kaḍugu}
  \item \emph{Say it:} \emph{kaḍugu}, once more
  \item \emph{Say it:} say \emph{kaḍugu}
  \item \emph{Recall:} say the Telugu for to call, then the Telugu for to count, then say \emph{kaḍugu} again
\end{itemize}

\begin{tcolorbox}[breakable,colback=teal!4,colframe=teal!35!black,title={Before you move on}]
What does \te{కడుగు} mean? (\textquotedblleft{}To wash\textquotedblright{}.) Say it once more.
\end{tcolorbox}
